% INSERTION: r1.pdf, p. 12, after the paragraph giving the four decay constants
% and the theta=90-degree cross-check (ending "reported in [34].").
% Place before "Because the initial and final mesons in each channel..."
% Upload figures/decay_constants_vs_theta.pdf to Prism and adjust the image
% path below if it is uploaded into a different directory.

The mixing-angle dependence of the decay constants is shown in
Fig.~\ref{fig:decay-constants-theta} for both heavy-strange sectors.
For each angle, $f_1(\theta)$ and $f_2(\theta)$ are recalculated from the
two-point sum rules, while the underlying QCD inputs, Borel parameters,
and continuum thresholds are held fixed. We use the same representative
input choices as in the angular scan of the radiative widths.
Within this scan, $f_1$ varies relatively weakly, whereas the variation of
$f_2$ is more pronounced, particularly in the bottom-strange sector.
These fixed-input curves illustrate the angular dependence and are distinct
from the propagated estimates quoted above.

\begin{figure}[tbp]
 \centering
 \includegraphics[width=\linewidth]{decay_constants_vs_theta.pdf}
 \caption{Mixing-angle dependence of the decay constants $f_1(\theta)$
 and $f_2(\theta)$ in the charm-strange (left) and bottom-strange (right)
 sectors. The solid blue curves correspond to the lower states
 $D_{s1}(2460)$ and the expected $B_{s1}(5750)$, while the dashed black
 curves correspond to $D_{s1}(2536)$ and $B_{s1}(5830)$.
 The red dash-dotted lines indicate the adopted central angles
 $\theta_{D_s}=26.6^\circ$ and $\theta_{B_s}=38.5^\circ$; the gray dotted
 lines mark the heavy-quark-limit reference value $35.3^\circ$.
 Both decay constants are recomputed at each angle with the remaining
 inputs fixed to the representative choices used in the radiative-width
 angular scan. The curves illustrate the angle dependence; no uncertainty
 bands are implied.}
 \label{fig:decay-constants-theta}
\end{figure}
